Mapa průměrných mezd ve zpracovatelském průmyslu (\ac{NACE} C) v~EU27 ukazuje, že ČR navzdory své silné průmyslové základně --- výrobní odvětví tvoří přes \SI{25}{\percent} \ac{HDP}, nejvíce v~\ac{EU} --- dosahuje mzdové úrovně výrazně pod svými produktivitními protihodnotami v~AT nebo DE.
Mzdová mezera v~průmyslu je emblematickým příkladem fenoménu, pro nějž se v~literatuře zakotvil pojem „low-wage manufacturing periphery“: výrobní kapacity přemístěné ze Západu do ČR přinesly zaměstnanost, ale ne přesun mzdové úrovně.
Tento pattern je udržitelný jen do doby, dokud výrobní náklady v~ČR zůstávají relativně příznivé --- s~rostoucí produktivitou, odlivem pracovní síly a~konkurencí o~pracovníky (srov.~obr.~\ref{fig:emigration_cz_timeline}) se diferenciál snižuje a~zaměstnavatelé budou nuceni ke mzdové konvergenci.
Proaktivní přístup --- tj.~sektorové \ac{KS} s~pravidelnou indexací průmyslových tarifů --- je z~pohledu pracovníků preferabilní oproti čekání na přirozený tržní tlak, který může trvat dekády a~probíhat za cenu dalšího odlivu pracovní síly.

Mapa průměrných mezd ve zpracovatelském průmyslu (\ac{NACE} C) v~EU27 ukazuje, že ČR navzdory své silné průmyslové základně --- výrobní odvětví tvoří přes \SI{25}{\percent} \ac{HDP}, nejvíce v~\ac{EU} --- dosahuje mzdové úrovně výrazně pod svými produktivitními protihodnotami v~AT nebo DE.
Mzdová mezera v~průmyslu je emblematickým příkladem fenoménu, pro nějž se v~literatuře zakotvil pojem „low-wage manufacturing periphery“: výrobní kapacity přemístěné ze Západu do ČR přinesly zaměstnanost, ale ne přesun mzdové úrovně.
Tento pattern je udržitelný jen do doby, dokud výrobní náklady v~ČR zůstávají relativně příznivé --- s~rostoucí produktivitou, odlivem pracovní síly a~konkurencí o~pracovníky (srov.~obr.~\ref{fig:emigration_cz_timeline}) se diferenciál snižuje a~zaměstnavatelé budou nuceni ke mzdové konvergenci.
Proaktivní přístup --- tj.~sektorové \ac{KS} s~pravidelnou indexací průmyslových tarifů --- je z~pohledu pracovníků preferabilní oproti čekání na přirozený tržní tlak, který může trvat dekády a~probíhat za cenu dalšího odlivu pracovní síly.

Mapa průměrných mezd ve zpracovatelském průmyslu (\ac{NACE} C) v~EU27 ukazuje, že ČR navzdory své silné průmyslové základně --- výrobní odvětví tvoří přes \SI{25}{\percent} \ac{HDP}, nejvíce v~\ac{EU} --- dosahuje mzdové úrovně výrazně pod svými produktivitními protihodnotami v~AT nebo DE.
Mzdová mezera v~průmyslu je emblematickým příkladem fenoménu, pro nějž se v~literatuře zakotvil pojem „low-wage manufacturing periphery“: výrobní kapacity přemístěné ze Západu do ČR přinesly zaměstnanost, ale ne přesun mzdové úrovně.
Tento pattern je udržitelný jen do doby, dokud výrobní náklady v~ČR zůstávají relativně příznivé --- s~rostoucí produktivitou, odlivem pracovní síly a~konkurencí o~pracovníky (srov.~obr.~\ref{fig:emigration_cz_timeline}) se diferenciál snižuje a~zaměstnavatelé budou nuceni ke mzdové konvergenci.
Proaktivní přístup --- tj.~sektorové \ac{KS} s~pravidelnou indexací průmyslových tarifů --- je z~pohledu pracovníků preferabilní oproti čekání na přirozený tržní tlak, který může trvat dekády a~probíhat za cenu dalšího odlivu pracovní síly.

Mapa průměrných mezd ve zpracovatelském průmyslu (\ac{NACE} C) v~EU27 ukazuje, že ČR navzdory své silné průmyslové základně --- výrobní odvětví tvoří přes \SI{25}{\percent} \ac{HDP}, nejvíce v~\ac{EU} --- dosahuje mzdové úrovně výrazně pod svými produktivitními protihodnotami v~AT nebo DE.
Mzdová mezera v~průmyslu je emblematickým příkladem fenoménu, pro nějž se v~literatuře zakotvil pojem „low-wage manufacturing periphery“: výrobní kapacity přemístěné ze Západu do ČR přinesly zaměstnanost, ale ne přesun mzdové úrovně.
Tento pattern je udržitelný jen do doby, dokud výrobní náklady v~ČR zůstávají relativně příznivé --- s~rostoucí produktivitou, odlivem pracovní síly a~konkurencí o~pracovníky (srov.~obr.~\ref{fig:emigration_cz_timeline}) se diferenciál snižuje a~zaměstnavatelé budou nuceni ke mzdové konvergenci.
Proaktivní přístup --- tj.~sektorové \ac{KS} s~pravidelnou indexací průmyslových tarifů --- je z~pohledu pracovníků preferabilní oproti čekání na přirozený tržní tlak, který může trvat dekády a~probíhat za cenu dalšího odlivu pracovní síly.

\input{texparts/python/sector_wages_map_C}



